% 12. Instrukcja uruchomienia
\section{Instrukcja uruchomienia}\label{sec:uruchomienie}

\begin{instrukcja}[wymagane]
Krok po kroku, tak żeby prowadząca mogła uruchomić projekt na swoim komputerze: co zainstalować, jakie
polecenia wydać, gdzie otworzyć aplikację. Sprawdźcie instrukcję na komputerze osoby, która nie
pisała kodu.
\end{instrukcja}

\begin{ramka}{Szybki start}
\textbf{Repozytorium:} \dkpokaz{repozytorium}\\
\textbf{Dane testowe:} aplikacja startuje z 12 przykładowymi terminami; konto nie jest potrzebne. Żeby wrócić do danych startowych, wystarczy usunąć klucz \texttt{planer-sesji:terminy} z \texttt{localStorage}.
\end{ramka}

\subsection{Wymagania}
\begin{itemize}
  \item Node.js 22.23.4 lub nowszy (z npm), Git
  \item przeglądarka: Chrome, Firefox albo Edge
\end{itemize}

\subsection{Uruchomienie lokalne}
\noindent\begin{minipage}{\linewidth}
\begin{lstlisting}[style=konsola]
git clone https://github.com/Machux133/MOIU-Planer-Sesji-I-Terminow.git
cd MOIU-Planer-Sesji-I-Terminow
npm install
npm run dev
# adres: http://localhost:4321/MOIU-Planer-Sesji-I-Terminow/
\end{lstlisting}
\end{minipage}
